\documentclass[aps,prx,preprint,superscriptaddress,nofootinbib,longbibliography]{revtex4-2}
\usepackage[T1]{fontenc}\usepackage{amsmath,amssymb,mathtools,amsthm,bm}\usepackage{booktabs,microtype,hyperref}
\newtheorem{theorem}{Theorem}\newtheorem{proposition}[theorem]{Proposition}\newtheorem{lemma}[theorem]{Lemma}\newtheorem{corollary}[theorem]{Corollary}\newtheorem{definition}[theorem]{Definition}
\newcommand{\E}{\mathbb E}\newcommand{\Prob}{\mathbb P}
\begin{document}
\title{Supplemental Material for ``Correction Exposure and Service Geometry in Sparse Hybrid Quantum Repeater Networks''}
\author{Ayush Nadiger}\affiliation{Department of Electrical and Computer Engineering, University of Massachusetts Amherst, Amherst, Massachusetts 01003, USA}
\date{August 31, 2026}\maketitle

\section{Odd-repetition transfer law and strict convexity}
Let $n=2r+1$ and let $q=(1-x)/2$ be the physical bit-flip probability written in the bias coordinate $x\in[0,1]$. Majority decoding fails when at least $r+1$ of the $2r+1$ bits flip. Define the logical bias $\lambda_n(x)=1-2q_{L,n}(x)$ and the exposure variable $u=-\log x$.

\begin{theorem}[Exact odd-repetition transfer]
With
\begin{equation}
K_r=\frac{(2r+1)\binom{2r}{r}}{4^r},
\end{equation}
one has
\begin{equation}
\lambda_{2r+1}(x)=K_r\int_0^x(1-t^2)^r\,dt.
\end{equation}
For
\begin{equation}
c_{2r+1}(u)=-\log\lambda_{2r+1}(e^{-u}),
\end{equation}
$c_{2r+1}''(u)>0$ for every $u>0$.
\end{theorem}
\begin{proof}
Differentiate the majority-vote binomial tail with respect to $x$. All noncentral terms telescope, leaving
\begin{equation}
\lambda'_{2r+1}(x)=K_r(1-x^2)^r,
\end{equation}
with $\lambda_{2r+1}(0)=0$, which yields the integral formula. Set
\begin{equation}
f(x)=(1-x^2)^r,\qquad F(x)=\int_0^x f(t)dt.
\end{equation}
The normalization $K_r$ cancels from logarithmic derivatives. With $x=e^{-u}$,
\begin{equation}
c_n'(u)=h(x),\qquad h(x)=\frac{x f(x)}{F(x)},\qquad c_n''(u)=-x h'(x).
\end{equation}
It is enough to show $h'(x)<0$ for $0<x<1$. Direct differentiation reduces the sign to positivity of
\begin{equation}
G(x)=x(1-x^2)^{r+1}-\bigl[1-(2r+1)x^2\bigr]F(x).
\end{equation}
Now $G(0)=0$ and
\begin{equation}
G'(x)=2x\bigl[(2r+1)F(x)-x f(x)\bigr].
\end{equation}
Because $f$ is strictly decreasing on $(0,1)$,
\begin{equation}
F(x)=\int_0^x f(t)dt>x f(x).
\end{equation}
Hence $(2r+1)F(x)-xf(x)>0$, so $G'(x)>0$, $G(x)>0$, $h'(x)<0$, and $c_n''(u)>0$.
\end{proof}

For REP3,
\begin{align}
\lambda_3(e^{-u})&=\frac{3e^{-u}-e^{-3u}}{2},\\
c_3(u)&=u+\log2-\log(3-e^{-2u}),\\
c_3''(u)&=\frac{12e^{2u}}{(3e^{2u}-1)^2}>0.
\end{align}
The small- and large-exposure limits are
\begin{align}
c_3(u)&=\frac32u^2+O(u^3),\\
c_3(u)&=u-\log(3/2)+o(1).
\end{align}

\section{Majorization, midpoint insertion, and correction charge}
Strict convexity gives Jensen's inequality, and Karamata strengthens it to the full majorization order.

\begin{corollary}[Strict exposure majorization]
For exposure vectors $\bm u$ and $\bm v$ with equal dimension and equal total exposure, if $\bm u$ majorizes $\bm v$, then
\begin{equation}
\sum_i c_n(u_i)\ge\sum_i c_n(v_i),
\end{equation}
with equality only for a permutation when all entries lie in the strictly convex domain.
\end{corollary}

For one gap of total exposure $U$, inserting a checkpoint at exposure coordinate $x\in(0,U)$ changes protected damage by
\begin{equation}
\Delta_n(U,x;\eta)=c_n(U)-c_n(x)-c_n(U-x)-\eta,
\end{equation}
where $\eta$ is the declared checkpoint charge in log-bias-damage units.

\begin{proposition}[Midpoint insertion]
For fixed $U$ and $\eta$, the unique maximizer of $\Delta_n(U,x;\eta)$ is $x=U/2$.
\end{proposition}
\begin{proof}
Differentiate:
\begin{equation}
\partial_x\Delta_n=-c_n'(x)+c_n'(U-x).
\end{equation}
Strictly increasing $c_n'$ gives a unique zero at $x=U/2$, with positive derivative to the left and negative derivative to the right.
\end{proof}

Define the midpoint splitting gain
\begin{equation}
S_n(U)=c_n(U)-2c_n(U/2).
\end{equation}
Strict convexity gives
\begin{equation}
S_n'(U)=c_n'(U)-c_n'(U/2)>0\qquad(U>0),
\end{equation}
so $S_n$ rises strictly from $S_n(0)=0$. Since $\lambda_{2r+1}(x)=K_rx+o(x)$ as $x\to0$, one has $c_n(U)=U-\log K_r+o(1)$ as $U\to\infty$, and therefore
\begin{equation}
\lim_{U\to\infty}S_n(U)=\eta_c(n)=\log K_r.
\end{equation}
Thus zero-cost midpoint insertion strictly improves every positive gap; for $0<\eta<\eta_c(n)$ there is a unique finite activation threshold $U_\eta$ and insertion is favorable exactly for $U>U_\eta$; for $\eta\ge\eta_c(n)$ no finite gap gives strict net benefit. For REP3, $\eta_c=\log(3/2)$.

For a homogeneous total exposure $U$ split into $m$ equal corrected segments,
\begin{equation}
D_m(U)=m c_n(U/m)+(m-1)\eta.
\end{equation}
Treating $m$ continuously and writing $u=U/m$, stationarity gives
\begin{equation}
u c_n'(u)-c_n(u)=\eta.
\end{equation}
The left-hand side has derivative $u c_n''(u)>0$, rises from zero, and approaches $\log K_r$, producing the same three correction-charge regimes.

\section{Hardware elasticity}
For memory exposure $u=\tau/T_2$ define
\begin{equation}
\Lambda_n(u)=\frac{u c_n'(u)}{c_n(u)}.
\end{equation}
Then
\begin{equation}
\frac{\partial\log c_n(\tau/T_2)}{\partial\log T_2}=-\Lambda_n(u).
\end{equation}
Since a length-$2r+1$ majority decoder first fails at order $q^{r+1}$ and $q=u/2+O(u^2)$,
\begin{equation}
c_n(u)=\Theta(u^{r+1})
\end{equation}
as $u\to0$, implying $\Lambda_n(u)\to r+1$. At high exposure the logical channel returns to first-order bias decay and $\Lambda_n(u)\to1$. The elasticity is a code-side local sensitivity, not a direct secret-key-rate derivative.

\section{Selective capability monotonicity}
Let $H$ be the installed hybrid-capable set and let $\Pi(H)$ denote the feasible set of routes and activation decisions under Selective semantics.

\begin{theorem}[Optional capability is monotone]
If $H\subseteq H'$ and unused installed capability has zero objective charge, then for any maximized service objective $J$,
\begin{equation}
J^\star(H')\ge J^\star(H).
\end{equation}
\end{theorem}
\begin{proof}
Every policy feasible under $H$ remains feasible under $H'$ by never activating sites in $H'\setminus H$. Thus $\Pi(H)\subseteq\Pi(H')$, and maximization over the larger feasible set cannot reduce the optimum.
\end{proof}
This theorem is a semantic audit. If a simulator reports worse \emph{optimized} performance after adding free optional capability, the route/activation constraints or optimizer must be checked before interpreting the result physically.

\section{Exact fixed-bundle three-EGC readiness}
For a fixed three-rail bundle, let $d_e\in\{1,2,3\}$ be the number of rails using geometric edge $e$. Each edge exposes three independent elementary-generation channels with per-slot success probability $q$. By slot $t$, one channel has succeeded with probability
\begin{equation}
x_t=1-(1-q)^t.
\end{equation}
Define
\begin{equation}
B_{3,d}(x)=\sum_{j=d}^{3}\binom3j x^j(1-x)^{3-j}.
\end{equation}

\begin{theorem}[Three-EGC readiness]
For independent edge processes,
\begin{equation}
\Prob(T_{\rm EGC}\le t)=\prod_{e\in E_{\rm bundle}}B_{3,d_e}(x_t),
\end{equation}
and
\begin{equation}
\E T_{\rm EGC}=\sum_{t\ge0}\left[1-\prod_eB_{3,d_e}(x_t)\right].
\end{equation}
\end{theorem}
\begin{proof}
The number of successful local channels on edge $e$ by time $t$ is $\mathrm{Binomial}(3,x_t)$. Edge $e$ can support its bundle multiplicity exactly when this count is at least $d_e$, giving $B_{3,d_e}(x_t)$. Physical-edge independence gives the product CDF. The mean follows from the tail-sum identity for nonnegative integer-valued completion time.
\end{proof}

The REP3 unfolding inequalities are explicit polynomials:
\begin{align}
B_{3,1}(x)^2-B_{3,2}(x)&=x^2(1-x)^2(x^2-4x+6)\ge0,\\
B_{3,2}(x)B_{3,1}(x)-B_{3,3}(x)&=x^3(1-x)(2x^2-7x+8)\ge0.
\end{align}
The residual quadratics are positive on $x\in[0,1]$. Therefore moving multiplicity $3$ from one edge to independent demands $2+1$, and then to $1+1+1$, improves the readiness CDF pointwise when geometric stretch is accounted for separately.

\section{Cut bottlenecks}
Every complete REP3 segment crossing a cut with $c$ usable independent edges must allocate total multiplicity three across at most $\min(3,c)$ edges. Repeatedly balancing any allocation $(u,v)$ with $u\ge v+2$ increases the product of binomial-tail factors by the log-concavity of binomial upper tails in their required multiplicity. Hence the best cut-only readiness CDF is
\begin{equation}
U_{3,c}(x)=
\begin{cases}
B_{3,3}(x),&c=1,\\
B_{3,2}(x)B_{3,1}(x),&c=2,\\
B_{3,1}(x)^3,&c\ge3.
\end{cases}
\end{equation}
A bridge therefore imposes an irreducible three-rail synchronization tax. This is a cut-only lower bound on completion time, not an SKR bound because detour and downstream timing remain outside it.

\section{Mandatory route signatures}
The fixed-bundle theorem conditions on a chosen three-rail route bundle. A separate, controller-dependent model applies when candidate physical rails arrive before the common checkpoint skeleton is committed.

\begin{definition}[Mandatory hybrid signature]
For a physical path $P$ and deployed hybrid set $H$, $\sigma_H(P)$ is the ordered sequence of mandatory hybrids encountered by $P$, terminals included. If $P$ is sampled from a declared path distribution, let $p_\sigma=\Prob[\sigma_H(P)=\sigma]$.
\end{definition}

\begin{proposition}[Route-signature collision]
For $q$ independent candidate rails,
\begin{equation}
\Phi_q(H)=\Prob(\text{all }q\text{ signatures agree})=\sum_\sigma p_\sigma^q.
\end{equation}
Refining the signature partition by adding a mandatory hybrid cannot increase $\Phi_q$.
\end{proposition}
\begin{proof}
Condition on the common signature. If a class of mass $p$ is refined into child masses $p_1,\ldots,p_m$ with $\sum_jp_j=p$, then $\sum_jp_j^q\le(\sum_jp_j)^q=p^q$ for integer $q\ge1$.
\end{proof}

\begin{theorem}[First compatible block under marked-Poisson arrivals]
Suppose candidate rails arrive as a Poisson process of total rate $\lambda$ and receive independent signature marks with probabilities $p_\sigma$. Let $T_q$ be the first time any signature class accumulates $q$ rails. Then
\begin{equation}
\Prob(T_q>t)=e^{-\lambda t}\prod_\sigma\left[\sum_{j=0}^{q-1}\frac{(\lambda p_\sigma t)^j}{j!}\right].
\end{equation}
Refining the signature partition delays $T_q$ pathwise under the natural coupling.
\end{theorem}
\begin{proof}
Poisson thinning makes class counts independent Poisson random variables of means $\lambda p_\sigma t$. The event $T_q>t$ is exactly that every class count is less than $q$, giving the product. Under partition refinement each coarse count is the sum of its child counts; whenever a child reaches $q$, its parent has already reached at least $q$, proving pathwise delay.
\end{proof}
For $m$ equally likely signatures,
\begin{equation}
\frac{\lambda T_q}{m^{(q-1)/q}}\Rightarrow X_q,
\qquad \Prob(X_q>x)=e^{-x^q/q!},
\end{equation}
so
\begin{equation}
\E T_q\sim\frac{(q!)^{1/q}\Gamma(1+1/q)}{\lambda}\Phi_q^{-1/q}.
\end{equation}
For REP3 the coefficient is $6^{1/3}\Gamma(4/3)\simeq1.62265$. The complete finite-time law depends on the full mass spectrum $\{p_\sigma\}$; $\Phi_3$ is an exact compatibility moment, not a sufficient state variable.

\section{Matched-gap separation witness}
To show that correction geometry and signature geometry are independent summaries, enumerate the 70 monotone shortest paths from $(0,0)$ to $(4,4)$ on the $5\times5$ square grid. There exist two four-hybrid deployments for which the complete route-level correction-gap histogram is identical: 35 paths have gap vector $(1,2,5)$ and 35 have $(1,7)$ in both deployments. Yet the signature-class masses differ. One deployment has ranked class counts $(20,20,15,10,5)$ and the other $(25,20,10,10,5)$, yielding
\begin{align}
\Phi_3^{(A)}&=0.05977, & \E T_3^{(A)}&=6.087/\lambda,\\
\Phi_3^{(B)}&=0.07507, & \E T_3^{(B)}&=5.855/\lambda.
\end{align}
Thus no statistic of correction gaps alone determines opportunistic three-rail reconvergence.

\section{Simulation provenance and evidence separation}
The manuscript intentionally uses multiple numerical generations for different claims.
\begin{enumerate}
\item The \textbf{registered reduced model} is a theorem-facing fidelity model calibrated to the Haldar--Guha--Towsley--Rozp\k{e}dek Type-1/Type-2 REP3 architecture. At the central point it uses $p_\ell=0.01$, $T_2^{(1)}=500$ nearest-neighbor heralding intervals, a tenfold Type-2 coherence improvement, Type-1 two-qubit error $10^{-3}$, a tenfold Type-2 gate improvement, zero interface error, and zero added local-operation duration. The 1,122 matched Forced/Selective instances are a semantic null control inside this generation.
\item The \textbf{event-driven density campaign} includes mandatory encoded actions and renewal timing and is used only to exhibit placement-dependent density regimes. Its absolute rate is not pooled with the reduced model.
\item The \textbf{small mechanism controls} isolate reconvergence/readiness changes at matched local conditions; they support the sign of the service mechanism, not a universal network optimum.
\item The \textbf{route-signature model} is a complementary opportunistic-assembly abstraction, not part of the fixed-bundle simulator. It is used for exact controller-dependent reconvergence laws and is kept separate from the reported event-driven rate values.
\end{enumerate}
The evidence hierarchy is therefore theorem $\rightarrow$ controlled mechanism $\rightarrow$ simulator-specific finite behavior. No simulator-specific density maximum is elevated into a topology-independent theorem.

\begin{thebibliography}{10}
\bibitem{Haldar2025} S. Haldar, S. Guha, D. Towsley, and F. Rozp\k{e}dek, Hybrid Repeaters with Encoding for Long Distance Entanglement Distribution, in \emph{2025 IEEE International Conference on Quantum Computing and Engineering (QCE)}, 1141--1147 (2025), doi:10.1109/QCE65121.2025.00128.
\end{thebibliography}
\end{document}
